\documentclass[11pt]{article}

\usepackage[T1]{fontenc}
\usepackage{lmodern}
\usepackage[margin=1in]{geometry}
\usepackage{amsmath,amssymb,bm}
\usepackage{booktabs,tabularx}
\usepackage{xcolor}
\usepackage{hyperref}

\setlength{\emergencystretch}{2em}

\hypersetup{
  colorlinks=true,
  linkcolor=blue!55!black,
  citecolor=blue!55!black,
  urlcolor=blue!55!black,
  pdftitle={Noncritical Type 0B String Theory and Its Matrix-Model Dual},
  pdfauthor={StringMC project}
}

\newcommand{\ap}{\alpha'}
\newcommand{\NSNS}{\mathrm{NS\!\!-\!\!NS}}
\newcommand{\RR}{\mathrm{R\!\!-\!\!R}}
\newcommand{\SL}{\mathrm{SL}}
\newcommand{\MQM}{\mathrm{MQM}}
\newcommand{\ellB}{\ell_{\mathrm B}}

\title{Noncritical Type 0B String Theory\\
and Its Matrix-Model Dual\\[0.4em]
\large Setup, conventions, and first benchmarks}
\author{StringMC project}
\date{\today}

\begin{document}

\maketitle

\begin{abstract}
This note specifies a new target for the StringMC project: the duality between
two-dimensional noncritical type 0B string theory and double-scaled Hermitian
matrix quantum mechanics with a symmetrically filled two-sided Fermi sea.  We
adopt explicitly the conventions of Balthazar--Rodriguez--Yin (BRY),
including $\ap=2$, their delta-normalized physical states, and their exact
matrix-model energy and coupling dictionary.  We identify the genus-one
logarithm as the first vacuum normalization benchmark and separate carefully
the sphere data supplied by BRY from the additional supermoduli measure
needed at genus one and genus two.
\end{abstract}

\tableofcontents

\section{Target and baseline}

The target is the proposed exact duality
\begin{equation}
 \boxed{
 \text{2D noncritical type 0B string theory}
 \quad\longleftrightarrow\quad
 \text{two-sided double-scaled MQM}}
 \label{eq:duality}
\end{equation}
developed in Refs.~\cite{DouglasEtAl2003NewHat,TakayanagiToumbas2003Type0B}.
Our master convention is Ref.~\cite{BalthazarRodriguezYin2023SMatrix}.
The essential distinction from the bosonic $c=1$ matrix model is the vacuum:
both sides of an even potential are filled to the same Fermi level.  The
resulting state has a nonperturbatively stable completion and supports two
collective modes, which map to the two propagating type 0B fields.

The initial comparison is defined by
\begin{equation}
 \begin{aligned}
 b&=1,
 &Q&=b+b^{-1}=2,
 &\ellB&\equiv\sqrt{\ap/2},
 &[\ellB]_{\rm BRY}&=1,\\
 q_{\RR}&=\widetilde q_{\RR}=0,
 &E_F&=-\mu_{\rm F},
 &\mu_{\rm F}&>0.
 \end{aligned}
 \label{eq:baseline}
\end{equation}
with equal regulated fillings $N_L=N_R$.  For the finite-radius benchmark,
Euclidean time is compactified as
\begin{equation}
 X\sim X+2\pi R_{\rm phys}.
 \label{eq:circle}
\end{equation}
The BRY notation $\ap=2$ is a choice of units:
$\ap=2\ellB^2$ followed by $[\ellB]_{\rm BRY}=1$.  It is not an equality
between a dimensionful quantity and the pure number two.
Every formula imported from another source is translated to this
$\ap=2$ convention before use.

\section{Worldsheet description}

\subsection{\texorpdfstring{$\widehat c=1$}{c-hat=1} matter and
super-Liouville theory}

The matter multiplet is the Lorentzian free field
$(X^0,\psi^0,\widetilde\psi^0)$.  In the BRY component normalization, the
$N=1$ super-Liouville action is
\begin{align}
 S_{\SL}=\int d^2z\bigg[
 &\frac{1}{2\pi}\left(
 \partial\phi\,\bar\partial\phi+
 \psi\,\bar\partial\psi+
 \widetilde\psi\,\partial\widetilde\psi\right)\notag\\
 &+2i\mu_{\rm L} b^2\psi\widetilde\psi e^{b\phi}
 +2\pi b^2\mu_{\rm L}^2e^{2b\phi}\bigg].
 \label{eq:sl-action}
\end{align}
The linear-dilaton slope and ordinary central charge are
\begin{equation}
 Q=b+b^{-1},
 \qquad
 c_{\SL}=\frac32(1+2Q^2).
 \label{eq:sl-central-charge}
\end{equation}
At $b=1$, $Q=2$ and $c_{\SL}=27/2$; together with the free timelike
multiplet this gives an $N=1$, $c=15$ matter SCFT.

We reserve $\phi$ for the Liouville boson and
$\varphi,\widetilde\varphi$ for the bosonized superghosts.  Throughout,
$h$ and $\widetilde h$ denote the holomorphic and antiholomorphic conformal
weights.  In the $\ap=2$ convention the elementary holomorphic weights are
\begin{equation}
 h\!\left[e^{\pm i\omega X^0}\right]=-\frac{\omega^2}{2},
 \qquad
 h\!\left[e^{q\varphi}\right]=-\frac12q(q+2),
 \qquad
 h[c]=-1,
 \qquad
 h[\sigma^0]=h[\mu^0]=\frac1{16},
 \label{eq:elementary-conformal-weights}
\end{equation}
with identical formulas in the antiholomorphic sector.  Thus
$h[e^{-\varphi}]=1/2$ and $h[e^{-\varphi/2}]=3/8$.  Equation
\eqref{eq:elementary-conformal-weights} fixes in one place the timelike,
superghost, and spin-field conventions used below.

The super-Liouville wall cuts off the strong-coupling end of the linear-dilaton
direction.  The weak-coupling asymptotic end supplies the scattering region.
BRY print the glyph $\mu$ both for the Liouville interaction in
Eq.~\eqref{eq:sl-action} and for the matrix-model Fermi depth.  We denote
them by $\mu_{\rm L}$ and $\mu_{\rm F}$ respectively.  They are not
identified inside the genus-one derivation; the comparison map
$\mu_{\rm L}=\kappa\mu_{\rm F}^{\,p}$ is introduced only afterward.

\subsection{Type 0B projection and spectrum}

The nonchiral type 0B GSO projection retains
\begin{equation}
 (\mathrm{NS}-,\mathrm{NS}-)
 \oplus(\mathrm{NS}+,\mathrm{NS}+)
 \oplus(\mathrm{R}+,\mathrm{R}+)
 \oplus(\mathrm{R}-,\mathrm{R}-).
 \label{eq:gso}
\end{equation}
There are no spacetime fermions.  The propagating fields are an $\NSNS$
scalar $T$ and an $\RR$ axion $A$.  The field traditionally called the
``tachyon'' is massless in the two-dimensional linear-dilaton background.
Discrete states also occur at special momenta, but they are not part of the
first vacuum-amplitude benchmark.

For $P\ge0$, BRY normalize the NS and R super-Liouville primaries by
\begin{align}
 h_{\rm NS}(P)&=\frac12\left(\frac{Q^2}{4}+P^2\right)
               =\frac{1+P^2}{2},\notag\\
 h_{\rm R}(P)&=\frac1{16}
 +\frac12\left(\frac{Q^2}{4}+P^2\right)
 =\frac1{16}+\frac{1+P^2}{2},
 \label{eq:liouville-weights}
\end{align}
Here $h_{\rm R}$ is the weight of the \emph{Liouville} Ramond operator;
the timelike spin or disorder field contributes a further $1/16$.  Hence
the combined matter Ramond operator used in the physical axion vertex has
\begin{equation}
 h[\mathcal A_P^\pm]=\widetilde h[\mathcal A_P^\pm]
 =\frac1{16}+h_{\rm R}(P)
 =\frac18+\frac12(1+P^2).
 \label{eq:combined-ramond-weight}
\end{equation}

The reflection phases are fixed by the weak-coupling, $\phi\to-\infty$,
asymptotics
\begin{align}
 V_P&\sim S_{\rm NS}(P)^{-1/2}e^{(Q/2+iP)\phi}
       +S_{\rm NS}(P)^{1/2}e^{(Q/2-iP)\phi},
 \label{eq:ns-reflection-asymptotic}\\
 V_P^{\mathrm R,+}&\sim\sigma^1\!\left[
 S_{\rm R}(P)^{-1/2}e^{(Q/2+iP)\phi}
 +S_{\rm R}(P)^{1/2}e^{(Q/2-iP)\phi}\right],
 \label{eq:rplus-reflection-asymptotic}\\
 V_P^{\mathrm R,-}&\sim\mu^1\!\left[
 S_{\rm R}(P)^{-1/2}e^{(Q/2+iP)\phi}
 -S_{\rm R}(P)^{1/2}e^{(Q/2-iP)\phi}\right].
 \label{eq:rminus-reflection-asymptotic}
\end{align}
Here $\sigma^1$ and $\mu^1$ are the spin and disorder fields of the
asymptotic Liouville Majorana fermion.  The relative minus sign in
Eq.~\eqref{eq:rminus-reflection-asymptotic} is part of the Ramond phase
convention, not an in/out label.  The explicit functions $S_{\rm NS}$ and
$S_{\rm R}$ are not needed in the BRY sphere amplitudes.

With those phases, the delta-function normalization is
\begin{equation}
 \langle V_{P_1}(z,\bar z)V_{P_2}(0)\rangle
 =\frac{\pi\delta(P_1-P_2)}
 {|z|^{2(h_1+h_2)}}.
 \label{eq:liouville-two-point}
\end{equation}
For either Ramond family $\epsilon=\pm$, the corresponding normalization is
\begin{equation}
 \left\langle
 V_{P_1}^{\mathrm R,\epsilon}(z,\bar z)
 V_{P_2}^{\mathrm R,\epsilon}(0)
 \right\rangle
 =\frac{\pi\delta(P_1-P_2)}
 {|z|^{2(h^{\rm R}_1+h^{\rm R}_2)}}.
 \label{eq:ramond-liouville-two-point}
\end{equation}
In the modular-invariant bosonic-CFT description,
$V_P^{\mathrm R,+}$ is a local lowest-weight operator, whereas
$V_P^{\mathrm R,-}$ is a defect operator attached to the $\mathbb Z_2$
fermion-parity topological line.  Thus the two symbols do not denote two
ordinary local-primary families.  The equal-family Ramond correlators below
contain either zero or two defect endpoints; a mixed $+$/$-$ correlator is
not an ordinary local sphere correlator in the BRY data.
Equations~\eqref{eq:liouville-two-point} and
\eqref{eq:ramond-liouville-two-point}, together with the standard free-field,
spin-field, and ghost inner products, are what is meant below by the
delta-normalized asymptotic string states.  They should not be confused with
a two-punctured sphere vacuum diagram, which retains a conformal-Killing
volume.
For the timelike zero mode we use the Fourier convention
\begin{equation}
 \int\frac{dX_0^0}{2\pi}
 e^{i(\omega-\omega_1-\omega_2)X_0^0}
 =\delta(\omega-\omega_1-\omega_2).
 \label{eq:timelike-zero-mode-normalization}
\end{equation}
Accordingly, the energy delta function below carries no additional $2\pi$;
that factor is absorbed into the zero-mode path-integral measure.

\subsection{Physical vertex operators, pictures, and sphere three-point data}

The incoming and outgoing physical operators are
\begin{align}
 T_\omega^\pm
 &=g_s c\widetilde c\,
 e^{-\varphi-\widetilde\varphi}
 e^{\pm i\omega X^0}V_{P=\omega},
 \label{eq:tachyon-vertex}\\
 A_\omega^\pm
 &=g_s\omega c\widetilde c\,
 e^{-\varphi/2-\widetilde\varphi/2}
 e^{\pm i\omega X^0}\mathcal A_{P=\omega}^{\pm},
 \label{eq:axion-vertex}\\
 \mathcal A_P^\pm
 &=\frac1{\sqrt2}\left(
 \sigma^0V_P^{\mathrm R,+}
 \pm\mu^0V_P^{\mathrm R,-}\right).
 \label{eq:ramond-combination}
\end{align}
Equivalently, their weak-coupling wave functions are
\begin{align}
 \mathcal A_P^+
 &\sim \mathsf S^+\widetilde{\mathsf S}^{+}S_{\rm R}(P)^{-1/2}
 e^{(Q/2+iP)\phi}
 +\mathsf S^-\widetilde{\mathsf S}^{-}S_{\rm R}(P)^{1/2}
 e^{(Q/2-iP)\phi},\notag\\
 \mathcal A_P^-
 &\sim \mathsf S^-\widetilde{\mathsf S}^{-}S_{\rm R}(P)^{-1/2}
 e^{(Q/2+iP)\phi}
 +\mathsf S^+\widetilde{\mathsf S}^{+}S_{\rm R}(P)^{1/2}
 e^{(Q/2-iP)\phi},
 \label{eq:combined-ramond-asymptotic}
\end{align}
where $\mathsf S^\pm$ and $\widetilde{\mathsf S}^{\pm}$ are the spin fields
of the combined timelike--Liouville fermion system, normalized so that
\begin{equation}
 \mathsf S^\pm\widetilde{\mathsf S}^{\pm}
 =\frac1{\sqrt2}\left(\sigma^0\sigma^1\pm\mu^0\mu^1\right).
 \label{eq:combined-spin-fields}
\end{equation}
Here $T_\omega^\pm$ is an $\NSNS$ vertex in picture $(-1,-1)$ and
$A_\omega^\pm$ is an $\RR$ vertex in picture
$(-\tfrac12,-\tfrac12)$.  For $\omega>0$, the superscript has the fixed
scattering interpretation
\begin{equation}
 +\;=\;\text{incoming},
 \qquad
 -\;=\;\text{outgoing}.
 \label{eq:in-out-sign}
\end{equation}
Thus an incoming state carries $e^{+i\omega X^0}$ and an outgoing state
carries $e^{-i\omega X^0}$ in these conventions.  The symbols $\pm$ on the
Liouville operators $V_P^{\mathrm R,\pm}$ instead label the two Ramond
families in Eq.~\eqref{eq:gso}; they are not in/out labels.  The combination
in Eq.~\eqref{eq:ramond-combination} correlates these two distinct uses of a
sign.  BRY choose their phases so that
Eqs.~\eqref{eq:liouville-two-point} and
\eqref{eq:ramond-liouville-two-point} hold and all structure constants below
are real for $P_i\geq0$.  No additional leg-pole factor is included.

The BRST mass-shell condition is $P=\omega$.  Before ghost dressing, the
left-moving matter weights are
\begin{align}
 h_T^{\rm matter}(P,\omega)
 &=-\frac{\omega^2}{2}+h_{\rm NS}(P)
 =\frac12+\frac{P^2-\omega^2}{2},\notag\\
 h_A^{\rm matter}(P,\omega)
 &=-\frac{\omega^2}{2}+h[\mathcal A_P^\pm]
 =\frac58+\frac{P^2-\omega^2}{2}.
 \label{eq:physical-matter-weight-audit}
\end{align}
At $P=\omega$, adding $e^{-\varphi}$ to the NS operator or
$e^{-\varphi/2}$ to the R operator gives weight one; multiplication by
$c$ then gives weight zero.  The same check holds on the right, so the
unintegrated physical vertices in Eqs.~\eqref{eq:tachyon-vertex} and
\eqref{eq:axion-vertex} have total weight $(0,0)$.

To state the three-point data independently of the ghost dressing, define
the level-$1/2$ descendants and the top NS component by
\begin{equation}
 \Lambda_P=G_{-1/2}V_P,
 \qquad
 \widetilde\Lambda_P=\widetilde G_{-1/2}V_P,
 \qquad
 W_P=G_{-1/2}\widetilde G_{-1/2}V_P.
 \label{eq:top-ns-component}
\end{equation}
BRY normalize the holomorphic supercurrent modes by
\begin{equation}
 \{G_r,G_s\}=2L_{r+s}
 +\frac{c}{12}(4r^2-1)\delta_{r,-s},
 \label{eq:sca-normalization}
\end{equation}
and analogously in the antiholomorphic sector.
The top component has two-point function
\begin{equation}
 \langle W_{P_1}(z,\bar z)W_{P_2}(0)\rangle
 =\frac{\pi(2h_1)^2\delta(P_1-P_2)}
 {|z|^{2(h_1+h_2+1)}}.
 \label{eq:top-ns-two-point}
\end{equation}
Writing $z_{ij}=z_i-z_j$, the two independent NS three-point functions are
\begin{align}
 \left\langle V_{P_1}(z_1)V_{P_2}(z_2)V_{P_3}(z_3)\right\rangle
 &=\frac{C(P_1,P_2,P_3)}
 {|z_{12}|^{2(h_1+h_2-h_3)}
  |z_{13}|^{2(h_1+h_3-h_2)}
  |z_{23}|^{2(h_2+h_3-h_1)}},
 \label{eq:ns-primary-three-point}\\
\left\langle W_{P_1}(z_1)V_{P_2}(z_2)V_{P_3}(z_3)\right\rangle
&=\frac{\widetilde C(P_1,P_2,P_3)}
 {|z_{12}|^{2(h_1+h_2-h_3+1/2)}
  |z_{13}|^{2(h_1+h_3-h_2+1/2)}
 |z_{23}|^{2(h_2+h_3-h_1-1/2)}}.
 \label{eq:ns-top-three-point}
\end{align}
\paragraph{Convention boundary for all Type~0B calculations.}
Equations~\eqref{eq:top-ns-two-point} and
\eqref{eq:ns-top-three-point} define the BRY-native coefficient
$\widetilde C_{\rm BRY}$.  It is real for positive continuum momenta and is
used directly, with no extra factor of $i$, in BRY-native sphere correlators,
PCO formulas, and the $G,H,J$ decompositions below.  This is different from
importing the same data into the graded sewing basis of the separate Human
Note.  If $D$ denotes the matched primary two-point normalization, the raw
ordered Human-Note state
\begin{equation}
 w_{\rm HN}=(G_{-1/2}V)\otimes
 (\widetilde G_{-1/2}\widetilde V)
 \label{eq:hn-raw-top-state}
\end{equation}
has
\begin{equation}
 \langle w_{\rm HN},w_{\rm HN}\rangle_{\rm HN}
 =-D(2h)^2,
 \qquad
 \langle W_{\rm BRY},W_{\rm BRY}\rangle_{\rm BRY}
 =+D(2h)^2.
 \label{eq:bry-hn-top-bpz-comparison}
\end{equation}
The minus sign follows directly from the Human Note's mutually
anticommuting left/right odd operators and its graded tensor-product BPZ
rule.  Consequently the pants dictionary at that interface is
\begin{equation}
 C_{\rm HN}^{(0)}=C_{\rm BRY},\qquad
 C_{\rm HN}^{(1)}=\sigma i\,\widetilde C_{\rm BRY},
 \qquad \sigma=\pm1,
 \label{eq:bry-hn-pants-dictionary}
\end{equation}
and the convention-independent content is
$(C_{\rm HN}^{(1)})^2=-\widetilde C_{\rm BRY}^{2}$.  The code takes
$\sigma=+1$ as a square-root branch.  A two-pants genus-two vacuum amplitude
cannot distinguish it from $\sigma=-1$; fixing that sign for an odd number
of odd pants requires an explicitly ordered component matrix element.  The
conversion is made exactly once at a BRY-to-Human-Note sewing boundary.  It
does not modify the Human Note's descendant signs and is unrelated to both
the physical free Majorana denominator and the auxiliary double-Virasoro
Majorana.  Numerical genus-two agreement checks this direct BPZ dictionary;
it is not used to infer it.

For two Ramond fields and one NS field, let $h_i^{\rm R}$ denote the
Ramond weights.  The nonzero equal-family correlators are
\begin{equation}
 \left\langle
 V_{P_3}(z_3)V_{P_2}^{\mathrm R,\epsilon}(z_2)
 V_{P_1}^{\mathrm R,\epsilon}(z_1)
 \right\rangle
 =\frac{C_\epsilon(P_1,P_2;P_3)}
 {|z_{12}|^{2(h_1^{\rm R}+h_2^{\rm R}-h_3)}
  |z_{13}|^{2(h_1^{\rm R}+h_3-h_2^{\rm R})}
  |z_{23}|^{2(h_2^{\rm R}+h_3-h_1^{\rm R})}},
 \qquad \epsilon=\pm,
 \label{eq:rrns-three-point}
\end{equation}
where
\begin{equation}
 C_\pm=\frac12\left(C_{\rm even}\pm C_{\rm odd}\right).
 \label{eq:cpm-definition-early}
\end{equation}
The coefficients are completely explicit at $b=1$.  Define
\begin{align}
 \Upsilon_{\rm NS}(x)
 &=\frac{\Gamma(x/2)}{\Gamma(1-x/2)}
   \Upsilon_1(x/2)^2,
 &
 \Upsilon_{\rm R}(x)
 &=\Upsilon_1\left(\frac{x+1}{2}\right)^2,\notag\\
 N_{\rm NS}(P)
 &=\frac{\Gamma(1+iP)}{\Gamma(1-iP)}
   \Upsilon_{\rm NS}(2iP),
 &
 N_{\rm R}(P)
 &=\frac{\Gamma(\frac12+iP)}{\Gamma(\frac12-iP)}
   \Upsilon_{\rm R}(2iP),
 \label{eq:upsilon-definitions}
\end{align}
The $b=1$ Barnes-double-gamma convention is fixed, including its overall
normalization, by \cite{BalthazarRodriguezYin2019LongString}
\begin{align}
 G(1)&=1,
 &G(x+1)&=\Gamma(x)G(x),\notag\\
 \Gamma_1(x)&=\frac{(2\pi)^{(x-1)/2}}{G(x)},
 &\Gamma_1(x+1)&=\frac{\sqrt{2\pi}}{\Gamma(x)}\Gamma_1(x),\notag\\
 \Upsilon_1(x)&=\frac{1}{\Gamma_1(x)\Gamma_1(2-x)},
 &\Upsilon_1(1)&=1.
 \label{eq:upsilon-one-definition}
\end{align}
Here $G$ is the Barnes $G$-function on its standard analytic continuation.
Equivalently, $\Gamma_1$ is meromorphic with poles at the nonpositive
integers, while $\Upsilon_1$ is entire, has zeros at
$x\in\mathbb Z\setminus\{1\}$, and obeys
\begin{equation}
 \Upsilon_1(2-x)=\Upsilon_1(x),
 \qquad
 \Upsilon_1(x+1)=\frac{\Gamma(x)}{\Gamma(1-x)}\Upsilon_1(x).
 \label{eq:upsilon-one-identities}
\end{equation}
The power of $2\pi$ in Eq.~\eqref{eq:upsilon-one-definition} uses the real
$\log(2\pi)$, so it introduces no additional complex-log branch choice.
At $P=0$, and whenever individual Gamma factors appear singular, the
composite expressions in Eqs.~\eqref{eq:upsilon-definitions}--\eqref{eq:c-odd}
are defined by analytic continuation and their continuous $P\to0^+$ limit,
not by separately evaluating the singular factors.

With
\begin{equation}
 P_\Sigma=P_1+P_2+P_3,
 \qquad
 \Delta_i=P_j+P_k-P_i,
 \qquad \{i,j,k\}=\{1,2,3\},
 \label{eq:three-point-momentum-combinations}
\end{equation}
the two NS coefficients are
\begin{align}
 \boxed{
 C(P_1,P_2,P_3)
 =\frac{i}{2}
 \frac{\prod_{j=1}^3N_{\rm NS}(P_j)}
 {\Upsilon_{\rm NS}(1+iP_\Sigma)
  \prod_{i=1}^3\Upsilon_{\rm NS}(1+i\Delta_i)}} ,
 \label{eq:c-ns}\\
 \boxed{
 \widetilde C(P_1,P_2,P_3)
 =i
 \frac{\prod_{j=1}^3N_{\rm NS}(P_j)}
 {\Upsilon_{\rm R}(1+iP_\Sigma)
  \prod_{i=1}^3\Upsilon_{\rm R}(1+i\Delta_i)}} .
 \label{eq:ctilde-ns}
\end{align}
For $P_1,P_2$ Ramond and $P_3$ NS, set
$\mathcal N=N_{\rm R}(P_1)N_{\rm R}(P_2)N_{\rm NS}(P_3)$.  Then
\begin{align}
 \boxed{
 C_{\rm even}(P_1,P_2;P_3)
 =-\frac{i}{2}\frac{\mathcal N}{
 \Upsilon_{\rm R}(1+iP_\Sigma)
 \Upsilon_{\rm R}(1+i\Delta_3)
 \Upsilon_{\rm NS}(1+i\Delta_1)
 \Upsilon_{\rm NS}(1+i\Delta_2)}} ,
 \label{eq:c-even}\\
 \boxed{
 C_{\rm odd}(P_1,P_2;P_3)
 =-\frac{i}{2}\frac{\mathcal N}{
 \Upsilon_{\rm NS}(1+iP_\Sigma)
 \Upsilon_{\rm NS}(1+i\Delta_3)
 \Upsilon_{\rm R}(1+i\Delta_1)
 \Upsilon_{\rm R}(1+i\Delta_2)}} .
 \label{eq:c-odd}
\end{align}
These coefficients are real for $P_i\geq0$ in the BRY phase convention.
Under analytic continuation that reflects one momentum,
$C$ and $\widetilde C$ change sign; reflecting the NS momentum changes the
sign of both $C_{\rm even}$ and $C_{\rm odd}$, while reflecting either
Ramond momentum exchanges their even and odd labels.  All descendant
correlators are then fixed by superconformal Ward identities.

The holomorphic PCO used in their sphere calculation is
\begin{equation}
 \chi=-\frac12e^\varphi
 \left(-i\psi^0\partial X^0+G^{\rm L}\right)
 +\text{ghost terms},
 \label{eq:bry-pco}
\end{equation}
and similarly in the antiholomorphic sector.  At higher genus this local
operator represents odd-modulus integration only after a global
supermoduli or vertical-integration prescription has been supplied.
For the tree-level NS amplitudes, BRY place a holomorphic and an
antiholomorphic PCO on a tachyon insertion.  Suppressing the spectator
$c\widetilde c$ factor, the resulting picture-raised operator is
\begin{align}
 \widetilde\chi\chi\,T_\omega^\pm
 =\frac{ig_s}{4}\big[&
 \omega^2\psi^0\widetilde\psi^0e^{\pm i\omega X^0}V_\omega
 +e^{\pm i\omega X^0}W_\omega\notag\\
 &\ \pm\omega\widetilde\psi^0e^{\pm i\omega X^0}\Lambda_\omega
 \mp\omega\psi^0e^{\pm i\omega X^0}\widetilde\Lambda_\omega
 \big]+\text{ghost terms}.
 \label{eq:picture-raised-tachyon}
\end{align}
The explicit $i$ in Eq.~\eqref{eq:picture-raised-tachyon} is BRY's cocycle
phase: combining the separately bosonized left- and right-moving factors
$e^{-\varphi}e^{-\widetilde\varphi}$ into
$e^{-\varphi-\widetilde\varphi}$ produces this extra factor.  It is already
included in the amplitude convention below.
Their ordered integration over a complex-conjugate pair of Grassmann-odd
moduli inserts $-i\widetilde\chi\chi$.  This sign is included in the
sphere amplitude below.

For completeness, distinguish the CFT three-point functions above from the
on-shell three-string amplitude.  The total picture on the sphere is
$(-2,-2)$.  For one incoming and two outgoing right-side modes
$\mathcal R=T+A$, place the incoming vertex at infinity and the outgoing
vertices at $1$ and $0$.  In the BRY cocycle and odd-modulus convention,
with $C_{S^2}$ denoting the sphere-topology normalization,
\begin{align}
 S_{1\mathcal R\to2\mathcal R}(\omega,\omega_1,\omega_2)
 &=i\,\delta(\omega-\omega_1-\omega_2)g_s^3C_{S^2}
 \bigg[
 \frac14\widetilde C(\omega,\omega_1,\omega_2)
 +\frac12\omega_1\omega_2
 C_{\rm even}(\omega_1,\omega_2;\omega)\notag\\
 &\hspace{3.8cm}
 +\frac12\omega\omega_2
 C_{\rm odd}(\omega,\omega_2;\omega_1)
 +\frac12\omega\omega_1
 C_{\rm odd}(\omega,\omega_1;\omega_2)
 \bigg].
 \label{eq:physical-sphere-three-point}
\end{align}
On the energy-conserving locus $\omega=\omega_1+\omega_2$, the bracket
collapses to $\omega\omega_1\omega_2$, so after removing the delta function
\begin{equation}
 \boxed{
 \mathcal A_{1\mathcal R\to2\mathcal R}^{S^2}
 =i g_s^3C_{S^2}\omega\omega_1\omega_2
 =\frac{4i}{\mu_{\rm F}}\omega\omega_1\omega_2.}
 \label{eq:physical-sphere-three-point-onshell}
\end{equation}
The same formula holds for three left-side modes $\mathcal L=T-A$.
Every sphere three-point amplitude containing both an $\mathcal L$ and an
$\mathcal R$ mode vanishes.  In particular, in the $T,A$ basis the zero-flux
selection rule eliminates amplitudes with an odd number of axions.

The odd spin structure contributes to the type 0B torus amplitude.  It is
therefore essential to retain the complete spin-structure sum: using only the
three even structures gives the wrong coefficient for the two massless
spacetime fields.

The three-punctured-sphere ledger above and the four-punctured-sphere block
layer below are intentionally separate.  The former fixes external-state and
pants normalizations; the latter computes chiral NS superconformal blocks but
does not yet perform the modulus integral of a four-string amplitude.

\section{Sphere four-point superconformal blocks}
\label{sec:sphere-superconformal-blocks}

This section fixes the code convention for the $N=1$ super-Virasoro block on
a four-punctured sphere, following BRY and the $c$-recursion of
Refs.~\cite{HadaszJaskolskiSuchanek2007NSBlocks,Belavin2007FourPoint,
Suchanek2011EllipticRecursion,BalthazarRodriguezYin2023SMatrix}.
Only NS external and exchanged representations are included at this stage.
The pointwise NS structure-constant contraction and the integral over real
internal Liouville momentum are included.  Ramond blocks, the integral over
the sphere modulus and its counterterms, analytic continuation to physical
scattering kinematics, and every higher-genus sewing problem remain outside
the scope of the present implementation.

\subsection{Labels, weights, and local series}

Parametrize the super-Liouville central charge and weights by
\begin{equation}
 c=\frac32(1+2Q^2),\qquad Q=b+b^{-1},\qquad
 h_i=\frac12\left(\frac{Q^2}{4}-a_i^2\right),\qquad
 h=\frac12\left(\frac{Q^2}{4}+P^2\right).
 \label{eq:ns-block-weights}
\end{equation}
For the type 0B theory, $b=1$, $Q=2$, $c=27/2$, and an external
super-Liouville primary of momentum $P_i$ has $h_i=(1+P_i^2)/2$.
At punctures 2 and 3 the external state may be either a primary or its
level-$1/2$ descendant.  Write
\begin{equation}
 \star^0h_i\equiv h_i,\qquad \star^1h_i\equiv *h_i,
 \qquad \underline h_i\equiv h_i+\frac{\delta_i}{2},
 \qquad \delta_i\in\{0,1\},\quad i=2,3.
 \label{eq:ns-block-star-notation}
\end{equation}
Thus there are eight chiral blocks
$\mathcal F^{e/o}(h_4,\star^{\delta_3}h_3,
\star^{\delta_2}h_2,h_1;h|z)$, where $e/o$ selects integer/half-integer
levels in the internal NS module.  The cross ratio and our local series are
\begin{align}
 z&=\frac{z_{12}z_{34}}{z_{13}z_{24}},\notag\\
 \mathcal F^e(z)
 &=z^{h-h_1-\underline h_2}
 \left[1+\sum_{m=1}^{\infty}z^m\mathfrak F_m\right],\notag\\
 \mathcal F^o(z)
 &=z^{h-h_1-\underline h_2}
 \sum_{n=0}^{\infty}z^{n+1/2}\mathfrak F_{n+1/2}.
 \label{eq:ns-block-local-series}
\end{align}
The fraktur coefficients specify the chiral phase used in the code: they are
the signed coefficients produced by BRY's ancillary notebook, including the
minus signs in the half-integer seeds below, with no second overall sign
inserted in the last line of Eq.~\eqref{eq:ns-block-local-series}.  BRY's
displayed block convention instead places an extra minus sign in front of the
odd series when $(\delta_3,\delta_2)=(1,1)$.  This is a holomorphic phase
choice, which the scalar decomposition into $|\mathcal F|^2$ does not fix;
the code choice is made explicit because it reproduces BRY's executable
crossing benchmark.

\subsection{The \texorpdfstring{$c$}{c}-recursion}

For $r\geq2$, $s\geq1$, and $r+s$ even, define
\begin{align}
 b_{r,s}(h)^2
 &=-\frac{4h+rs-1+
 \sqrt{16h^2+8(rs-1)h+(r-s)^2}}{r^2-1},\notag\\
 c_{r,s}(h)&=\frac{15}{2}+3b_{r,s}(h)^2+3b_{r,s}(h)^{-2},\notag\\
 A_{r,s}(h)&=\frac12
 \prod_{\substack{p=1-r,\ldots,r\\q=1-s,\ldots,s\\
 (p,q)\ne(0,0),(r,s)\\p+q\in2\mathbb Z}}
 \frac{\sqrt2}{p b_{r,s}(h)+q b_{r,s}(h)^{-1}}.
 \label{eq:ns-block-pole-data}
\end{align}
With the external labels suppressed, the recursion is
\begin{align}
 \mathfrak F_0(h;c)&=1,\notag\\
 \mathfrak F_\ell(h;c)
 &=f_\ell(h)+
 \sum_{\substack{r\geq2,\ s\geq1\\r+s\in2\mathbb Z\\rs\leq2\ell}}
 \frac{R^{\ell}_{r,s}(h)}{c-c_{r,s}(h)}
 \mathfrak F_{\ell-rs/2}
 \left(h+\frac{rs}{2};c_{r,s}(h)\right).
 \label{eq:ns-block-c-recursion}
\end{align}
The endpoint $rs=2\ell$ is essential: it supplies the first $(3,1)$ pole
at $\ell=3/2$ and the first $(2,2)$ pole at $\ell=2$.  BRY's displayed
equation prints the strict inequality $rs<2\ell$, whereas its ancillary
notebook sums through $rs=2\ell$.  Equation~\eqref{eq:ns-block-c-recursion}
and the code follow the notebook.

At a pole set $b=b_{r,s}(h)$, $Q_{r,s}=b+b^{-1}$, and
\begin{equation}
 a_i^{(r,s)}=\sqrt{\frac{Q_{r,s}^2}{4}-2h_i}.
 \label{eq:ns-block-pole-momenta}
\end{equation}
For $\delta\in\{0,1\}$, define the fusion polynomial
\begin{align}
 P_{r,s}^{(\delta)}(h_a,h_b)
 ={}&\prod_{\substack{p=1-r,\,1-r+2,\ldots,r-1\\
 q=1-s,\,1-s+2,\ldots,s-1\\
 (p+q)-(r+s)\equiv2(1-\delta)\ ({\rm mod}\ 4)}}
 \frac{2a_a^{(r,s)}-2a_b^{(r,s)}-pb-q/b}{2\sqrt2}\notag\\
 &\hspace{2.2cm}\times
 \frac{2a_a^{(r,s)}+2a_b^{(r,s)}+pb+q/b}{2\sqrt2}.
 \label{eq:ns-block-fusion-polynomial}
\end{align}
Let $\delta_i^{(\ell)}=\delta_i$ for integer $\ell$ and
$\delta_i^{(\ell)}=1-\delta_i$ for half-integer $\ell$.  The residue is
\begin{equation}
 R^{\ell}_{r,s}(h)=
 (-1)^{rs\delta_3}
 \left(-\frac{\partial c_{r,s}(h)}{\partial h}\right)A_{r,s}(h)
 P_{r,s}^{(\delta_2^{(\ell)})}(h_1,h_2)
 P_{r,s}^{(\delta_3^{(\ell)})}(h_4,h_3).
 \label{eq:ns-block-residue}
\end{equation}
The factor $(-1)^{rs\delta_3}$ is required; omitting it is the sign typo in
the older recursion formula noted and corrected by BRY.

For an implementation-complete seed ledger, set
$A=h+h_3-h_4$, $B=h+h_2-h_1$.  For $m\in\mathbb Z_{>0}$ and
$n\in\mathbb Z_{\geq0}$,
\begin{equation}
\renewcommand{\arraystretch}{1.55}
\begin{array}{c|c|c}
 (\delta_3,\delta_2)&f_m&f_{n+1/2}\\ \hline
 (0,0)&\displaystyle\frac{(A)_m(B)_m}{m!(2h)_m}
 &\displaystyle\frac{(A+\frac12)_n(B+\frac12)_n}{n!(2h)_{n+1}}\\
 (0,1)&\displaystyle\frac{(A)_m(B+\frac12)_m}{m!(2h)_m}
 &\displaystyle\frac{(A+\frac12)_n(B)_{n+1}}{n!(2h)_{n+1}}\\
 (1,0)&\displaystyle\frac{(A+\frac12)_m(B)_m}{m!(2h)_m}
 &\displaystyle-\frac{(A)_{n+1}(B+\frac12)_n}{n!(2h)_{n+1}}\\
 (1,1)&\displaystyle\frac{(A+\frac12)_m(B+\frac12)_m}{m!(2h)_m}
 &\displaystyle-\frac{(A)_{n+1}(B)_{n+1}}{n!(2h)_{n+1}}
\end{array}
\label{eq:ns-block-seeds}
\end{equation}
where $(x)_n$ is the rising Pochhammer symbol.

\subsection{Elliptic representation and code interface}

The local $z$-series is converted to the elliptic nome
\begin{equation}
 q(z)=\exp\!\left[-\pi\frac{K(1-z)}{K(z)}\right],
 \qquad K(z)={}_2F_1\!\left(\frac12,\frac12;1;z\right),
 \qquad z=\frac{\theta_2(q)^4}{\theta_3(q)^4}.
 \label{eq:ns-block-nome}
\end{equation}
Both parities are written as
\begin{align}
 \mathcal F^{e/o}(z)={}&(16q)^{h-Q^2/8}
 z^{Q^2/8-h_1-\underline h_2}
 (1-z)^{Q^2/8-\underline h_2-\underline h_3}\notag\\
 &\times\theta_3(q)^{\frac32Q^2-
 4(h_1+\underline h_2+\underline h_3+h_4)}
 H^{e/o}(q).
 \label{eq:ns-block-elliptic-factorization}
\end{align}
Here $H^e$ has integer powers of $q$, whereas $H^o$ has powers
$q^{n+1/2}$.  The elliptic series converges rapidly throughout the direct
channel and is the representation used for numerical four-punctured-sphere
work.

The root-level Python module
\texttt{superconformal\_}\allowbreak\texttt{blocks.py} implements
Eqs.~\eqref{eq:ns-block-weights}--\eqref{eq:ns-block-elliptic-factorization}.
Its main class stores a level $\ell$ by the integer $2\ell$, exposes
individual recursion coefficients, and evaluates either the local series or
the elliptic series:
\begin{verbatim}
block = NSSphereFourPointBlock.from_liouville_momenta(
    p1=P1, p2=P2, p3=P3, p4=P4,
    internal_momentum=P, b=1, star2=True, star3=True)
Fe = block.elliptic_block(z, order=8, parity="even")
Fo = block.elliptic_block(z, order=8, parity="odd")
\end{verbatim}
All square roots, fractional powers, and elliptic functions use principal
branches.  Crucially, the $a_i^{(r,s)}$ are recomputed with $Q_{r,s}$ at
every recursion pole rather than with the physical background charge.  A
near-zero $c-c_{r,s}$ raises an explicit pole error instead of silently
returning an unstable value.  The order parameter counts retained
coefficients: $q^0$ through $q^{N-1}$ in the even block and $q^{1/2}$
through $q^{N-1/2}$ in the odd block.

As an external regression anchor, take BRY's ancillary-notebook values
\begin{equation}
 c=\frac{27}{2}+10^{-5},\quad
 (P_4,P_3,P_2,P_1)=\left(\frac35,\frac14,\frac13,\frac12\right),\quad
 h=\frac{11}{10},\quad (\delta_3,\delta_2)=(1,1).
 \label{eq:ns-block-bry-benchmark-input}
\end{equation}
The implementation reproduces
\begin{align}
 H^e(q)&=1+1.0998737373737373q
 +15.824538587885476q^2+21.09367085588721q^3+\cdots,\notag\\
 H^o(q)&=-1.7823926767676768q^{1/2}
 -6.738809221289842q^{3/2}
 -13.848138977969848q^{5/2}+\cdots.
 \label{eq:ns-block-bry-benchmark-series}
\end{align}
At $z=1/3+3i/5$, truncation through eight coefficients gives
\begin{align}
 \mathcal F^e&=1.3109708952736965
 +0.18058411587390863i,\notag\\
 \mathcal F^o&=-0.3051714889777302
 -0.4379494400515853i.
 \label{eq:ns-block-bry-benchmark-values}
\end{align}
The automated tests also compare the $z$- and $q$-representations for all
eight choices and verify explicitly that both endpoint poles are present.
No claim about a higher-genus analogue is made here; that literature and
sewing audit is deferred to a later stage.

\subsection{Nonchiral correlators and the BRY four-tachyon density}
\label{subsec:bry-sphere-correlators}

We now contract the chiral blocks with the $b=1$ pants coefficients defined
in Eqs.~\eqref{eq:ns-primary-three-point} and
\eqref{eq:ns-top-three-point}, using the explicit formulas
\eqref{eq:c-ns} and \eqref{eq:ctilde-ns}.  Place the operators $4,3,2,1$
at $\infty,1,z,0$, respectively, and abbreviate
\begin{equation}
 C_{12P}=C(P_1,P_2,P),\quad C_{34P}=C(P_3,P_4,P),\quad
 \widetilde C_{12P}=\widetilde C(P_1,P_2,P),\quad
 \widetilde C_{34P}=\widetilde C(P_3,P_4,P).
 \label{eq:sphere-four-point-coefficient-shorthand}
\end{equation}
Below, $\mathcal F^{e/o}$ has external labels $(h_4,h_3,h_2,h_1)$, whereas
$\mathcal F^{e/o}_{**}$ has $(h_4,*h_3,*h_2,h_1)$.  BRY's three
nonchiral functions are
\begin{align}
 G(4321|z)={}&\int_0^\infty\frac{dP}{\pi}\bigl[
 C_{12P}C_{34P}\mathcal F^e(z)\mathcal F^e(\bar z)
 +\widetilde C_{12P}\widetilde C_{34P}
 \mathcal F^o(z)\mathcal F^o(\bar z)\bigr],
 \label{eq:bry-G-decomposition}\\
 H(4321|z)={}&-\int_0^\infty\frac{dP}{\pi}\bigl[
 \widetilde C_{12P}\widetilde C_{34P}
 \mathcal F^e_{**}(z)\mathcal F^e_{**}(\bar z)
 +C_{12P}C_{34P}\mathcal F^o_{**}(z)
 \mathcal F^o_{**}(\bar z)\bigr].
 \label{eq:bry-H-decomposition}
\end{align}
Here
$G=\langle V_{P_4}(\infty)V_{P_3}(1)V_{P_2}(z,\bar z)V_{P_1}(0)\rangle$
and
$H=\langle V_{P_4}(\infty)W_{P_3}(1)W_{P_2}(z,\bar z)V_{P_1}(0)\rangle$.
The third function includes the two timelike-free-fermion contractions in
the picture-raised tachyons:
\begin{align}
 J(4321|z)=-\int_0^\infty\frac{dP}{\pi}\Bigg\{&
 C_{12P}C_{34P}\left[
 \frac{\mathcal F^o_{**}(z)\mathcal F^e(\bar z)}{1-\bar z}
 +\frac{\mathcal F^e(z)\mathcal F^o_{**}(\bar z)}{1-z}\right]
 \notag\\[-2pt]
 &+\widetilde C_{12P}\widetilde C_{34P}\left[
 \frac{\mathcal F^e_{**}(z)\mathcal F^o(\bar z)}{1-\bar z}
 +\frac{\mathcal F^o(z)\mathcal F^e_{**}(\bar z)}{1-z}\right]
 \Bigg\}.
 \label{eq:bry-J-decomposition}
\end{align}
In the scalar convention used in Eqs.~\eqref{eq:bry-H-decomposition} and
\eqref{eq:bry-J-decomposition}, the doubly-starred odd block is minus the
signed ancillary-notebook series returned by the low-level block class.
This phase cancels between the two chiral halves of $H$ but is indispensable
in the mixed products of $J$.  For nonnegative real momenta, $H$ and $J$ are
real and negative, as required by the anti-Hermiticity of the descendant
operators.

For the all-tachyon term in BRY's $1\to3$ setup, assign
\begin{equation}
 (P_4,P_3,P_2,P_1)=(\omega,\omega_3,\omega_2,\omega_1).
 \label{eq:bry-four-tachyon-momentum-assignment}
\end{equation}
After the internal-$P$ integral but before the $z$ integral, the reduced
moduli density is
\begin{align}
 \mathcal I_T(z)={}&|z|^{-2\omega_1\omega_2}
 |1-z|^{-2\omega_2\omega_3}
 \left[\frac{\omega_2^2\omega_3^2}{|1-z|^2}G(4321|z)
 -H(4321|z)-\omega_2\omega_3J(4321|z)\right].
 \label{eq:bry-four-tachyon-density}
\end{align}
Thus the corresponding unregularized contribution carries the overall
factor
\begin{equation}
 \frac{i}{2}\,\delta\!\left(\omega-\sum_{i=1}^3\omega_i\right)
 g_s^4C_{S^2}\int_{\mathbb C}d^2z\,\mathcal I_T(z).
 \label{eq:bry-four-tachyon-unregularized-prefactor}
\end{equation}
This is the BRY expression after multiplying the explicitly calculated
four-tachyon term by the eight nonzero even-axion contributions.  Equality of
those eight terms follows from the assumed perturbative decoupling of left-
and right-side modes; it was not separately derived from eight worldsheet
correlators.

The code interface is
\begin{verbatim}
corr = BRYNSFourPointCorrelator(
    p1=1/2, p2=1/3, p3=1/4, p4=3/5, block_order=8)
values = corr.evaluate(z=0.1, p_max=5, quadrature_order=20)
\end{verbatim}
The module \texttt{sphere\_four\_point.py} performs the $P$ integral by
Gauss--Legendre quadrature, while
\texttt{super\_liouville\_structure\_constants.py} supplies the normalized
$C$ and $\widetilde C$.  The recursion is evaluated at
$c=27/2+10^{-5}$, as in BRY's numerical benchmark, and with arbitrary
precision internally.  The latter is necessary because separately large
near-threshold recursion terms cancel at higher orders.  The finite
$P_{\max}$ is only a numerical truncation: convergence must be checked by
increasing $P_{\max}$, the quadrature order, and the elliptic-block order.

At BRY's crossing-test momenta
$(P_1,P_2,P_3,P_4)=(1/2,1/3,1/4,3/5)$, using the parameters in the code
snippet, the regression values are
\begin{equation}
 G(0.1)=0.6251752343,\qquad
 H(0.1)=-0.3082787765,\qquad
 J(0.1)=-0.4125437927.
 \label{eq:bry-four-point-regression-values}
\end{equation}
With the same numbers interpreted as
$(\omega_1,\omega_2,\omega_3,\omega)$,
$\mathcal I_T(0.1)=0.7630630252$.  The automated tests check these anchors,
the $H$ crossing relation, structure-constant identities, and the
near-threshold high-precision cancellation.

\subsection{First regulated $1_R\to3_R$ benchmark}

The first end-to-end test follows BRY Eq.~(4.15):
\begin{equation}
 \omega=\frac13+ia,\qquad
 \omega_1=\omega_2=\omega_3=\frac{\omega}{3},\qquad
 0.5\leq a\leq0.7.
 \label{eq:bry-complex-energy-benchmark}
\end{equation}
No pole of the super-Liouville structure constants crosses the undeformed
contour $P\in\mathbb R_{\geq0}$ along this continuation.  The $s$- and
$u$-channel counterterms vanish, as does the t-channel term proportional to
$\widetilde C\widetilde C$.  Only the leading $C C$ subtraction at the
picture-changing collision $z=1$ remains.  At fixed internal momentum its
three picture-raised coefficients form the OPE square
\begin{align}
 B_t(P)
 &=\omega_2^2\omega_3^2+c_0(P)^2
   +2\omega_2\omega_3c_0(P)
 \notag\\
 &=\frac14\left[1+(\omega_2+\omega_3)^2-P^2\right]^2,
 &
 c_0(P)&=h_2+h_3-h_P.
 \label{eq:bry-leading-t-counterterm-coefficient}
\end{align}
Consequently the nonzero fixed-$P$ counterterm density, including $dP/\pi$,
is
\begin{equation}
 r_t(P,z)=\frac{1}{\pi}
 C(\omega_1,\omega,P)C(\omega_2,\omega_3,P)B_t(P)
 |1-z|^{-3+P^2-(\omega_2+\omega_3)^2},
 \quad 0\leq P<P_* ,
 \label{eq:bry-leading-t-counterterm-density}
\end{equation}
where
\begin{equation}
 P_*=\sqrt{1+\operatorname{Re}\!\left[(\omega_2+\omega_3)^2\right]}.
 \label{eq:bry-t-counterterm-threshold}
\end{equation}
For $a=0.6$, $P_*=0.9430708966$.  The subtraction is global, as required by
analytic continuation of the power divergence; it is not a literal removal
of only the $1/\epsilon$ term in a small disk.

Following BRY, inversion crossing folds the full $z$ plane to
$D=\{|z|\leq1\}$, which is divided into
\begin{equation}
 D_\epsilon^{(0)}=\{|z|<\epsilon\},\qquad
 D_{\rm int}=\{\epsilon\leq|z|\leq1, |1-z|\geq\epsilon\},\qquad
 D_\epsilon^{(1)}=\{|z|\leq1, |1-z|<\epsilon\}.
 \label{eq:bry-three-moduli-regions}
\end{equation}
We use $\epsilon=10^{-2}$.  The direct low-$z$ OPE is integrated in
$D_\epsilon^{(0)}$, the elliptic-recursion blocks are used in $D_{\rm int}$,
and the crossed low-$(1-z)$ blocks through next-to-leading order are used in
$D_\epsilon^{(1)}$.  The angular quadrature in $D_{\rm int}$ is explicitly
split at $\arg z=\pm\epsilon$; a single rule on $[-\pi,\pi]$ misses this
narrow boundary layer.

Let $\mathcal M$ denote the regularized expression inside braces in
BRY Eq.~(B.6), after the three $z$ integrals and the final $P$ integral.  The
worldsheet and matrix-model normalizations imply
\begin{align}
 \mathcal A_{1_R\to3_R}^{\rm WS}\big|_{\delta,\mu_{\rm F}^{-2}\ {\rm stripped}}
 &=\frac{8i}{\pi}\mathcal M,
 &
 \mathcal M_{\rm MQM}
 &=\pi\omega\omega_1\omega_2\omega_3(1+2i\omega),
 \label{eq:bry-reduced-moduli-target}
\end{align}
which is equivalent to
$\mathcal A_{1_R\to3_R}^{\rm MQM}
=8i\omega\omega_1\omega_2\omega_3(1+2i\omega)$.
The executable implementation is \texttt{bry\_one\_to\_three.py}.  It keeps
the internal-$P$ contour and threshold explicit, evaluates all three moduli
regions, and reports both sides of Eq.~\eqref{eq:bry-reduced-moduli-target}
as complex numbers.  Fast regression tests separately check the OPE square,
$P_*$, the direct low-$z$ expansion, complex external momenta, and the final
normalization conversion and regularized integration path.  At $a=0.6$, a
representative $q^8$ run with
$P_{\max}=4$, 24-point $P$ quadrature on each side of $P_*$, and the moduli
orders recorded by the command-line output gives
\begin{align}
 \mathcal M_{\rm WS}&=0.01750879-0.00320541i,
 &
 \mathcal M_{\rm MQM}&=0.01772691-0.00297257i,
 \notag\\
 \mathcal A_{\rm WS}&=0.00816252+0.04458577i,
 &
 \mathcal A_{\rm MQM}&=0.00756958+0.04514122i.
 \label{eq:bry-first-regulated-numerical-result}
\end{align}
The relative complex-amplitude discrepancy is $1.78\%$.  This recovers the
BRY curve at the first nontrivial level, but does not yet reproduce their
quoted sub-$0.3\%$ numerical precision.  At the same production quadrature,
the cutoff comparison begins with
\begin{align}
 \mathcal M_{q^6}&=0.01765123551-0.00248615555i,
 &
 \mathcal M_{q^8}&=0.01750879244-0.00320541444i,
 \notag\\
 \mathcal M_{q^{10}}&=0.01750781745-0.00322155122i.
 \label{eq:bry-first-regulated-q-ledger}
\end{align}
In units of $|\mathcal M_{\rm MQM}|$, the $q^6\to q^8$ displacement is
$4.08\%$, while the $q^8\to q^{10}$ displacement is approximately $0.09\%$;
changing the moduli orders from 12 to 14 at $q^8$ gives $0.074\%$.  Thus
$q^6$ is not converged and production begins at approximately $q^8$.  The
remaining limitation is the narrow-$z=1$ and oscillatory-$P$ product
quadrature rather than the superconformal recursion.

\section{Matrix quantum mechanics}

\subsection{Double-scaled Hamiltonian}

The BRY double-scaled matrix quantum mechanics is
\begin{equation}
 H=\frac12\operatorname{Tr}(P^2-X^2),
 \qquad
 h=\frac12(p^2-x^2).
 \label{eq:one-particle-hamiltonian}
\end{equation}
Diagonalization and the Vandermonde determinant turn the eigenvalues into
free fermions.  All one-particle states on both sides with
\begin{equation}
 E\le-\mu_{\rm F},\qquad \mu_{\rm F}>0,
 \label{eq:fermi-level}
\end{equation}
are filled.  No extra oscillator curvature, matrix Planck constant, or
$\ap$ is to be restored in Eq.~\eqref{eq:one-particle-hamiltonian}.

\subsection{Two-sided basis and exact dictionary}

BRY's physical left- and right-side scattering modes are
\begin{align}
 \mathcal L_\omega^\pm&=T_\omega^\pm-A_\omega^\pm,
 &
 \mathcal R_\omega^\pm&=T_\omega^\pm+A_\omega^\pm.
 \label{eq:lr-basis}
\end{align}
There is no $1/\sqrt2$ in this basis.  Reflection exchanges
$\mathcal L$ and $\mathcal R$ and realizes $(-1)^{F_L}$; hence $T$ is even
and $A$ is odd.  At zero flux, correlators with an odd number of axions
vanish.

The two sides decouple to every order in the perturbative
$1/\mu_{\rm F}$ expansion.
Eigenvalue tunnelling across the barrier mixes them through effects
nonanalytic in $g_s$.  Perturbative and tunnelling contributions must therefore
be tracked separately.

For one incoming and $k$ outgoing excitations on one side,
\begin{equation}
 \mathcal A_{1\to k}
 =\sum_{g\ge0}\mu_{\rm F}^{-(k-1+2g)}
 \mathcal A_{\rm pert}^{(g)}
 +O(e^{-2\pi\mu_{\rm F}}).
 \label{eq:mm-genus-expansion}
\end{equation}
The exact worldsheet/matrix-model dictionary is
\begin{equation}
 \boxed{
 \omega_{\rm MM}=2\omega,\qquad
 g_s=\frac4{\pi\mu_{\rm F}},\qquad
 C_{S^2}=\frac{\pi}{g_s^2}.}
 \label{eq:bry-dictionary}
\end{equation}
The energy relation includes the corresponding rescaling of the
energy-conservation delta function.  In the phases of
Eqs.~\eqref{eq:tachyon-vertex}--\eqref{eq:ramond-combination}, BRY compare
directly to the matrix-model S-matrix; an additional leg-pole transform
must not be inserted.

\section{Duality dictionary}

\begin{table}[ht]
\centering
\small
\begin{tabularx}{\textwidth}{@{}p{0.36\textwidth}X@{}}
\toprule
String description & Matrix-model description \\
\midrule
Liouville coordinate
& Eigenvalue time-of-flight coordinate \\
BRY Fermi parameter $\mu_{\rm F}$
& Depth of the Fermi level below the quadratic maximum \\
$g_s=4/(\pi\mu_{\rm F})$
& Expansion about the filled two-sided Fermi sea \\
$\NSNS$ scalar $T$
& Even combination of left/right scattering modes \\
$\RR$ axion $A$
& Odd combination of left/right scattering modes \\
Worldsheet energy $\omega$
& Matrix energy $\omega_{\rm MM}=2\omega$ \\
$(-1)^{F_L}$
& Reflection $M\mapsto-M$ \\
$\RR$ background data
& Fermi-level asymmetry and flux data \\
D0-brane excitation
& A fermion excited from the sea toward the top \\
Nonperturbative correction
& Eigenvalue tunnelling across the barrier \\
\bottomrule
\end{tabularx}
\caption{Baseline 0B/matrix-model dictionary.  Flux normalizations are
deferred until the flux milestone.}
\label{tab:dictionary}
\end{table}

Equation~\eqref{eq:lr-basis}, rather than an independently orthonormalized
density basis, is used in every amplitude comparison.

\section{Genus-one path integral in BRY variables}

We use the term free energy for the connected string vacuum functional,
$\mathcal F=\log Z_{\rm string}$; it is not the thermodynamic quantity
$-\beta^{-1}\log Z$.  This sentence fixes the overall sign convention.

Define the physical radius and its dimensionless BRY numerical value by
\begin{equation}
 \rho\equiv\frac{R_{\rm phys}}{\ellB},
 \qquad
 x\equiv\frac{R_{\rm phys}}{\sqrt{2\ap}}
 =\frac{\rho}{2},
 \qquad
 y\equiv\frac{R_{\rm phys}}{\sqrt{\ap}}
 =\frac{\rho}{\sqrt2}=\sqrt2x.
 \label{eq:dimensionless-radius}
\end{equation}
The variable $y$ is the argument used in the finite-radius formulas of
Ref.~\cite{DouglasEtAl2003NewHat}.  We denote the connected amplitude per
regulated BRY Liouville-coordinate length by
$\widehat{\mathcal F}_1=\mathcal F_1/V_\phi^{\rm reg}$.

\subsection{Even spin structures and the modular orbit integral}

We follow the worldsheet derivation of
Ref.~\cite{DouglasEtAl2003NewHat}, translated to Eq.~\eqref{eq:dimensionless-radius}.
For an even spin structure $(r,s)=(0,1),(1,0),(1,1)$, let $D_{r,s}$ be the
fermion determinant divided by the square root of the scalar determinant,
and define
\begin{equation}
 \Theta_\rho(\tau)=\sum_{m,n\in\mathbb Z}e^{-S_{m,n}},
 \qquad
 S_{m,n}=\frac{\pi\rho^2}{2\tau_2}|n-m\tau|^2.
 \label{eq:bry-circle-lattice}
\end{equation}
The compact matter multiplet, super-Liouville theory, and combined ghost
system contribute respectively
\begin{equation}
 Z^X_{r,s}=\frac{y}{\sqrt{\tau_2}}|D_{r,s}|^2\Theta_\rho,
 \quad
 Z^{\rm SL}_{r,s}=\frac{V_\phi^{\rm reg}}{2\pi\sqrt{2\tau_2}}
 |D_{r,s}|^2,
 \quad
 Z^{\rm gh}_{r,s}=\frac1{2\tau_2}|D_{r,s}|^{-4}.
 \label{eq:even-spin-factors}
\end{equation}
All nonzero-mode determinants cancel.  Including the diagonal GSO factor
$1/2$ for each of the three even structures gives
\begin{equation}
 \widehat{\mathcal F}_{1,\rm even}
 =\frac{3\rho}{16\pi}
 \int_{\mathcal F}\frac{d^2\tau}{\tau_2^2}\Theta_\rho(\tau).
 \label{eq:even-spin-integral}
\end{equation}

This modular integral can be evaluated without a divergent vacuum-energy
sum.  Separate the zero orbit and unfold every nonzero orbit to the strip:
\begin{align}
 I(\rho)&\equiv
 \int_{\mathcal F}\frac{d^2\tau}{\tau_2^2}\Theta_\rho(\tau)\notag\\
 &=\frac\pi3+2\sum_{k=1}^{\infty}
 \int_{-1/2}^{1/2}d\tau_1\int_0^\infty\frac{d\tau_2}{\tau_2^2}
 \exp\left[-\frac{\pi\rho^2k^2}{2\tau_2}\right]\notag\\
 &=\frac\pi3+\frac4{\pi\rho^2}\sum_{k=1}^{\infty}\frac1{k^2}
 =\frac\pi3\left(1+\frac2{\rho^2}\right).
 \label{eq:lattice-orbit-integral}
\end{align}
The series is absolutely convergent; the last equality uses the ordinary
Basel sum $\sum_{k\ge1}k^{-2}=\pi^2/6$, not analytic continuation at
$s=-1$.  Therefore
\begin{equation}
 \boxed{
 \widehat{\mathcal F}_{1,\rm even}
 =\frac\rho{16}+\frac1{8\rho}.}
 \label{eq:even-spin-density}
\end{equation}

\subsection{Odd spin structure and direct odd-supermodulus integration}

The odd structure $(r,s)=(0,0)$ has left- and right-moving $\gamma$ zero
modes from the conformal Killing spinors, $\beta$ zero modes from the
gravitino components that cannot be gauged away, and zero modes for both the
matter fermion $\chi$ and Liouville fermion $\psi$.  The Liouville-fermion
zeros cancel the $\gamma$-zero-mode divergence.  Integrating the $\beta$
zeros, equivalently the odd moduli, inserts the supercurrents
\begin{equation}
 G(z)\overline G(\bar w)
 =\big(\chi\partial X+\psi\partial\phi\big)(z)
  \big(\bar\chi\bar\partial X+\bar\psi\bar\partial\phi\big)(\bar w).
 \label{eq:odd-supercurrent-insertion}
\end{equation}
These insertions absorb the remaining $\chi,\bar\chi$ zero modes.  The
nontrivial matter factor is then
\begin{equation}
 \left\langle\partial X(z)\bar\partial X(\bar w)\right\rangle_X,
 \label{eq:odd-matter-correlator}
\end{equation}
including both its separated-point term $-\pi/\tau_2$ and its contact term.
Because the compact-boson action is proportional to
$\int\partial X\bar\partial X$, the integrated insertion is a radius
derivative of the circle path integral.  It fixes the odd contribution to
be proportional to $y-y^{-1}$.

Equivalently, one can work globally on the odd supertorus
\begin{equation}
 (z,\theta)\sim(z+1,\theta)
 \sim(z+\tau+\lambda\theta,\theta+\lambda)
 \label{eq:odd-supertorus}
\end{equation}
and integrate $d\lambda\,d\bar\lambda$ directly.  Thus no arbitrary local
PCO position is introduced.  If a local PCO description is used instead,
it must reproduce this direct integral and its contact term; there is no
additional vertical correction to add to a globally defined supermoduli
integral for this unpunctured genus-one vacuum amplitude.

The path integral fixes the radius dependence.  As in
Ref.~\cite{DouglasEtAl2003NewHat}, its overall magnitude is normalized by
the large-radius spectrum.  The result before choosing the 0A/0B sign is
\begin{equation}
 \widehat{\mathcal F}_{1,\rm odd}^{(\varepsilon)}
 =\frac{\varepsilon}{24\sqrt2}\left(y-\frac1y\right)
 =\varepsilon\left(\frac\rho{48}-\frac1{24\rho}\right).
 \label{eq:odd-spin-general}
\end{equation}
The sign is chosen so that the $1/R$ term counts one massless field in 0A
and two in 0B.  In this convention $\varepsilon_A=+1$ and
$\varepsilon_B=-1$, hence
\begin{equation}
 \boxed{
 \widehat{\mathcal F}_{1,\rm odd}^{0B}
 =-\frac\rho{48}+\frac1{24\rho}.}
 \label{eq:odd-spin-density}
\end{equation}

\subsection{The 0B and fluxless 0A spin sums}

The even contribution in Eq.~\eqref{eq:even-spin-density} is common to the
two theories.  Their relative GSO choice reverses the odd contribution, so
\begin{equation}
 \widehat{\mathcal F}_{1,\rm odd}^{0A}(\rho_A)
 =\frac{\rho_A}{48}-\frac1{24\rho_A}
 =-\widehat{\mathcal F}_{1,\rm odd}^{0B}(\rho_A).
 \label{eq:odd-spin-density-0a}
\end{equation}
Adding even and odd structures gives directly
\begin{align}
 \boxed{
 \widehat{\mathcal F}^{0B}_1(\rho_B)
 =\frac{\rho_B}{24}+\frac1{6\rho_B}}
 \label{eq:bry-spectrum-torus}\\[3pt]
 \boxed{
 \widehat{\mathcal F}^{0A}_1(\rho_A)
 =\frac1{12}\left(\rho_A+\frac1{\rho_A}\right)}.
 \label{eq:bry-spectrum-torus-0a}
\end{align}
The BRST spectrum explains these coefficients: 0B has two momentum towers
and one winding tower, while 0A has one momentum tower and two winding
towers.  This is now a check of the modular and supermoduli derivation, not
a zeta-regularized input.

The physical T-duality map is
\begin{equation}
 R_A=\frac{\ap}{R_B},
 \qquad \rho_A=\frac2{\rho_B},
 \qquad
 \widehat{\mathcal F}^{0A}_1(\rho_A)
 =\widehat{\mathcal F}^{0B}_1\!\left(\frac2{\rho_A}\right).
 \label{eq:0a-0b-t-duality}
\end{equation}
This also maps the 0B stationary point $\rho_B=2$ to the 0A stationary
point $\rho_A=1$, or $R_A=\sqrt{\ap/2}$.

\subsection{Liouville wall and the regulator ledger}

The interaction in Eq.~\eqref{eq:sl-action} is invariant under
\begin{equation}
 \phi\longrightarrow\phi+c,
 \qquad
 \mu_{\rm L}\longrightarrow e^{-bc}\mu_{\rm L}.
 \label{eq:liouville-shift}
\end{equation}
Consequently the wall position contains
$-b^{-1}\log\mu_{\rm L}$.  With a fixed coordinate cutoff in the
weak-coupling region, define
\begin{equation}
 V_\phi^{\rm reg}
 =-\frac1b\log\frac{\mu_{\rm L}}{\Lambda_{\rm L}}
 +V_{\rm scheme}.
 \label{eq:bry-liouville-volume}
\end{equation}
The additive term depends on the location and shape of the cutoff.  It is
not compared.  Combining Eqs.~\eqref{eq:bry-spectrum-torus} and
\eqref{eq:bry-liouville-volume} at $b=1$ gives
\begin{equation}
 \boxed{
 \mathcal F^{0B}_1(\rho,\mu_{\rm L})\big|_{\rm univ}
 =-\frac{\log(\mu_{\rm L}/\Lambda_{\rm L})}{24}
 \left(\rho+\frac4\rho\right)}.
 \label{eq:torus-anchor}
\end{equation}
For fluxless type 0A, the same regulated volume gives
\begin{equation}
 \boxed{
 \mathcal F^{0A}_1(\rho_A,\mu_{\rm L})\big|_{\rm univ}
 =-\frac1{12}\left(\rho_A+\frac1{\rho_A}\right)
 \log\frac{\mu_{\rm L}}{\Lambda_{\rm L}}.}
 \label{eq:torus-anchor-0a}
\end{equation}
Restoring dimensions, this is
\begin{equation}
 \mathcal F^{0A}_1\big|_{\rm univ}
 =-\frac{\log(\mu_{\rm L}/\Lambda_{\rm L})}{12\sqrt2}
 \left(\frac{2R_A}{\sqrt{\ap}}+\frac{\sqrt{\ap}}{R_A}\right).
 \label{eq:torus-anchor-0a-dimensional}
\end{equation}
The regulator-independent datum is the logarithmic derivative
\begin{equation}
 \boxed{
 \frac{\partial\mathcal F^{0B}_1}
 {\partial\log\mu_{\rm L}}
 =-\frac1{24}\left(\rho+\frac4\rho\right).}
 \label{eq:torus-log-derivative}
\end{equation}

If another source uses $\Phi=a\phi$ and a parameter $\lambda$ defined by
$\mu_{\rm L}=\kappa\lambda^p$, then
\begin{equation}
 V_\Phi^{\rm reg}\big|_{\log\lambda}
 =-\frac{ap}{b}\log\lambda.
 \label{eq:volume-translation}
\end{equation}
Thus $\kappa$ and additive cutoff changes affect only a constant, whereas a
field rescaling $a$ or power map $p$ changes the logarithmic coefficient.

\subsection{Douglas convention and the matrix variable}

The calculation above follows Sec.~4.2 and App.~A of
Ref.~\cite{DouglasEtAl2003NewHat}, with $y=R_{\rm phys}/\sqrt\ap$
translated to the BRY radius $\rho=\sqrt2y$.  That reference writes the
Liouville kinetic term as
$(8\pi)^{-1}(d\phi)^2$, which is the same coordinate normalization as
Eq.~\eqref{eq:sl-action}, and then sets
$V_L=-\log|\mu_{\rm D}|$, with an additive cutoff constant suppressed.
At $b=1$ this is Eq.~\eqref{eq:bry-liouville-volume} after a constant
rescaling of the interaction parameter.  Their Eqs.~(4.5), (4.7), and
(4.8) become respectively Eqs.~\eqref{eq:even-spin-density},
\eqref{eq:odd-spin-general}, and the two totals above.

To compare the logarithm, introduce rather than assume
\begin{equation}
 \mu_{\rm L}=\kappa\mu_{\rm F}^{\,p}.
 \label{eq:liouville-fermi-map}
\end{equation}
Equation~\eqref{eq:torus-anchor} then predicts
\begin{equation}
 \mathcal F^{0B}_1\big|_{\log\mu_{\rm F}}
 =-\frac p{24}\left(\rho+\frac4\rho\right)\log\mu_{\rm F}.
 \label{eq:torus-fermi-comparison}
\end{equation}
The conventional string/matrix-model identification corresponds to
$p=1$ and an unobservable constant $\kappa$.  Agreement therefore tests
$p=1$; it is not an input to the worldsheet moduli integral.

For the smoke-test choice
$\mu_{\rm L}/\Lambda_{\rm L}=10$, the stationary numerical radius is
$\rho=2$, equivalently $R_{\rm phys}=\sqrt{2\ap}$, and
\begin{equation}
 \frac{\partial\mathcal F^{0B}_1}
 {\partial\log\mu_{\rm L}}\bigg|_{\rho=2}
 =-\frac16,
 \qquad
 \mathcal F^{0B}_1(\rho=2)\big|_{\rm univ}
 =-\frac{\log10}{6}.
 \label{eq:numerical-anchor}
\end{equation}

\section{Benchmark ladder}

The calculations should be developed in the following order.

\subsection*{A. Analytic torus check}

First reproduce the determinant cancellation in
Eq.~\eqref{eq:even-spin-factors}, unfold the convergent modular integral
Eq.~\eqref{eq:lattice-orbit-integral}, and derive the odd-spin term by
integrating the odd supermodulus with its contact term.  Then apply the
0A/0B sign and the Liouville wall displacement.  The spectrum count and
matrix variable $\mu_{\rm F}$ are subsequent checks.

\subsection*{B. Inverted-oscillator scattering data}

Construct even and odd one-particle scattering states, compute their reflection
phases and regulated densities of states, and reconstruct the finite-radius
genus-one term from the two parity sectors.  Cutoff variation and increasing
numerical precision should leave the logarithmic coefficient stable.

\subsection*{C. Collective fields and sphere amplitudes}

Build the left and right time-of-flight fields, form the BRY combinations
\(\mathcal L=T-A\) and \(\mathcal R=T+A\), verify the zero-flux parity
selection rule, and reproduce the BRY-normalized $1\to2$ and $1\to3$
amplitudes with $\omega_{\rm MM}=2\omega$.  Mixed-side connected amplitudes
must vanish perturbatively.

\subsection*{D. Flux backgrounds and 0A/0B T-duality}

Introduce asymmetric Fermi data using explicitly named flux variables.  Keep
Lorentzian continuous-flux conventions distinct from Euclidean quantized
fluxes, and verify the exact finite-radius 0A/0B map
\cite{GukovTakayanagiToumbas2004Flux,MaldacenaSeiberg2005FluxVacua}.

\subsection*{E. Higher genus and StringMC}

Only after the lower-genus dictionary is fixed should the project choose a
super-Riemann-surface measure, specify the treatment of odd moduli, and adapt
the existing StringMC sampling infrastructure.  The bosonic moduli measure
must not be imported into this target implicitly.  Degeneration,
factorization, GSO phases, and superghost normalization must be audited before
a full higher-genus integral is attempted.

\section{Worldsheet literature for the genus-one and genus-two program}

The higher-genus target requires more than a bosonic moduli integral with a
spin-structure prefactor.  It needs the BRST-complete state space, the
$bc$--$\beta\gamma$ measure, the integration over odd moduli, and the
super-Liouville conformal data in every allowed NS and Ramond sewing channel.
A detailed source-to-formula map is maintained in
\texttt{literature\_genus1\_genus2.md}; the main conclusions are summarized
here.

Ref.~\cite{BalthazarRodriguezYin2023SMatrix} computes tree-level sphere
$1\to2$ and $1\to3$ amplitudes.  It does not compute the torus vacuum
amplitude or the genus-two free energy, and explicitly leaves higher-genus
blocks and loop amplitudes for future work.  We therefore use it as the
master convention and sphere-data source, not as a substitute for the
higher-genus measure.

\subsection{Genus-one measure source}

The generic propagating spectrum consists of the NS--NS tachyon and R--R axion,
supplemented by discrete BRST states at special momenta.  The physical 0B
spectrum and its circle momentum/winding projection are given in
Ref.~\cite{DouglasEtAl2003NewHat}; the systematic NS and R BRST cohomology is
given in Ref.~\cite{ItohOhta1992BRST}.

Section~5 follows the explicit genus-one measure of
Ref.~\cite{DouglasEtAl2003NewHat}, translated to the BRY ledger.  For the
three even torus spin structures, the matter, super-Liouville, and combined
superghost factors are
\begin{align}
 Z_X&=\frac{y}{\sqrt{\tau_2}}|D_{r,s}|^2
       \sum_{m,n}e^{-S_{m,n}},\notag\\
 Z_{\SL}&=\frac{V_{\rm D}}{2\pi\sqrt{2\tau_2}}|D_{r,s}|^2,
 &
 Z_{\rm sgh}&=\frac{1}{2\tau_2}|D_{r,s}|^{-4}.
 \label{eq:torus-determinants}
\end{align}
All nonzero-mode determinants cancel.  An even supertorus has moduli
dimension $1|0$, so there is no PCO insertion for the unpunctured vacuum
amplitude.  The quantity $V_{\rm D}$ is their regulated Liouville volume;
it is translated using Eq.~\eqref{eq:volume-translation}, rather than
identified by its symbol.

The odd spin structure must not be omitted.  It has one odd supermodulus,
one $\beta$ and one $\gamma$ zero mode in each chiral sector, and matter and
Liouville fermion zero modes.  Appendix A of
Ref.~\cite{DouglasEtAl2003NewHat} integrates the odd modulus directly.  The
$\beta$ zero modes insert $G\overline G$, the Liouville-fermion zero modes
cancel the $\gamma$-zero-mode divergence, and the remaining insertion is
proportional to $\langle\partial X\,\overline\partial X\rangle$, including
its contact term.  This gives the odd-spin contribution required for
Eq.~\eqref{eq:torus-anchor}.

For this vacuum calculation, direct integration over the global odd
supertorus modulus is the primary prescription.  Vertical integration is
needed only if the same integral is represented by local PCO sections, or if
punctures or off-shell states destroy the global description.  It must then
reproduce, rather than be added to, the direct-supermoduli result.

\subsection{Genus two: conformal data and measure}

For an unpunctured surface of genus $g\ge2$,
\begin{equation}
 \dim\mathfrak M_g=(3g-3\,|\,2g-2).
 \label{eq:supermoduli-dimension}
\end{equation}
Thus the genus-two chiral measure has dimension $3|2$: a PCO
representation uses two holomorphic and two antiholomorphic PCOs.  There are
ten even and six odd genus-two spin structures
\cite{Witten2012SuperRiemannSurfaces}.  The all-genus type-0 GSO projection
is a diagonal spin-structure sum.  Type 0B assigns no Arf sign, whereas 0A
weights a spin structure by $(-1)^{\operatorname{Arf}}$
\cite{KaidiEtAl2020TopologicalSuperconductors}.

The complete three-point coefficient is a family of coefficients.  The
NS sector has two independent structures,
\begin{equation}
 C(P_1,P_2,P_3),\qquad \widetilde C(P_1,P_2,P_3),
\end{equation}
and an R--R--NS pants vertex requires two further structures, which may be
taken as $C_{\rm even}$ and $C_{\rm odd}$.  They were bootstrapped in
Refs.~\cite{RashkovStanishkov1996ThreePoint,Poghossian1997StructureConstants}.
For numerical work at $b=1$, the preferred normalization is the explicit
$\Upsilon_{\rm NS/R}$ representation of
Ref.~\cite{BalthazarRodriguezYin2023SMatrix}, whose formulas and reflection
rules have already been fixed explicitly in
Eqs.~\eqref{eq:upsilon-definitions}--\eqref{eq:c-odd}.
They are part of the initial worldsheet convention, not data introduced only
at genus two.  In particular, PCO supercurrents couple to top components and
descendants, so the bottom coefficient $C$ alone is insufficient.

The structure constants must be combined with super-Virasoro sewing blocks.
NS elliptic recursions, including top-component blocks, are developed in
Refs.~\cite{BelavinEtAl2007Bootstrap,Belavin2007FourPoint,Suchanek2011EllipticRecursion};
Ramond blocks are treated in Ref.~\cite{HadaszJaskolskiSuchanek2008RamondBlocks}.
These are benchmarks for a plumbing-fixture sewing engine, not a ready-made
genus-two vacuum block.

For even spin structures, the natural starting point is the slice-independent
D'Hoker--Phong measure based on the super period matrix
\cite{DHokerPhong2002ChiralMeasure,DHokerPhong2002Lectures}.  Their formula
for a general $N=1$, $c=15$ matter SCFT contains both the matter-supercurrent
correlator and the finite-dimensional $b$- and $\beta$-ghost determinants.
The latter are essential: the apparent collision pole of the naive product
of two PCOs is cancelled by an opposite singularity in those determinants.
For the present problem, the flat-space matter partition function must be
replaced by the noncompact time $\times$ super-Liouville theory, with its
Liouville zero-mode regulator and degeneration boundary terms retained.

The standard explicit D'Hoker--Phong result is an even-spin construction.
The six odd spin structures are therefore a separate work package.  Their
fermion and superghost zero modes, contact terms, and separating and
nonseparating degenerations must be derived directly before translating the
answer into a PCO gauge.

The BRY coupling relation does fix the conversion of genus powers,
\begin{equation}
 \frac1{\mu_{\rm F}^2}=\frac{\pi^2}{16}g_s^2.
 \label{eq:genus-two-power-map}
\end{equation}
Consequently a matrix-model term $A_2/\mu_{\rm F}^2$ maps to
$(\pi^2A_2/16)g_s^2$.  This power map does not determine the absolute
unpunctured genus-two path-integral normalization.  That constant must be
fixed by degeneration and factorization to the BRY-normalized sphere data
and the torus anchor.

\subsection{Vertical integration at genus two}

If PCOs are used globally, spurious poles obstruct a single smooth PCO
section.  Sen and Witten prescribe safe local sections and vertical
interpolations between them
\cite{Sen2015OffShell,SenWitten2015PCO}.  With two PCOs, the construction
contains bulk cell integrals, codimension-one vertical segments on shared
faces, and codimension-two vertical cells at corners; the hierarchy then
terminates.  The choices must also respect plumbing degenerations.  This is
an alternative representation of the odd-moduli integral, not an additional
physical correction to a correct super-period-matrix calculation.

\subsection{Result of the literature audit}

The primary literature supplies the complete sphere super-Liouville
three-point data and a powerful even-spin genus-two gauge-fixing framework.
It does not appear to supply the finished diagonal type-0B genus-two vacuum
integrand for time $\times$ super-Liouville including all odd spin
structures.  Constructing and factorization-testing that assembly is the
main worldsheet research problem for the genus-two target.

\section{Deferred choices and safeguards}

The baseline intentionally leaves the following items open:
\begin{itemize}
 \item continuation to $\mu_{\rm F}<0$ through the top of the potential;
 \item unequal filling and the two independent 0B flux variables;
 \item Lorentzian versus Euclidean flux normalization;
 \item the complete 0A/0B T-duality dictionary;
 \item the nonperturbative tunnelling contour prescription;
 \item the higher-genus supermoduli integration cycle and measure.
\end{itemize}
An imbalance between the Fermi seas is not a harmless numerical regulator: it
changes the $\RR$ background.  Likewise, Eq.~\eqref{eq:bry-dictionary} fixes
the sphere coupling convention but does not by itself normalize the
unpunctured genus-two integration cycle.  Both facts should remain explicit
in code and data products.

\section{Immediate next step}

The next bounded calculation is to complete stage A at the integrand level:
derive Eq.~\eqref{eq:odd-spin-density} by direct integration of the odd
supertorus modulus, including its contact term, without importing its value
from Douglas et al.  The matrix-model phase-shift calculation then tests
Eq.~\eqref{eq:liouville-fermi-map} and in particular the exponent $p$.

\bibliographystyle{unsrt}
\bibliography{../References/references}

\end{document}
